% 本LaTeX模板的一般使用说明
\chapter{说明\footnote{章标题中脚注命令测试}}
Again，这是北方民族大学学术论文\LaTeX{}模板（\CTeX{}-Based）---\NMUThesis{}。

本\LaTeX{}模板为北民大研究生学术论文模板，适用于理工类学术硕士和专业硕士。本\LaTeX{}模板参考自院教字〔2003〕169号《北方民族大学研究生学位论文格式和要求》，基本覆盖了论文内容和格式方面的要求。如图\ref{fig:first_image_t}是北方民族大学校标。

\begin{figure}[htbp]
	\centering
	\includegraphics[width=.5\textwidth]{nmu_logo}
	\caption{北方民族大学校标}\label{fig:first_image_t}
\end{figure}
%-----------------------------
\section{宏包\footnote{节标题中脚注命令测试}使用}

请将以下文件与此LaTeX文件放在同一目录中\footnote{本模板改写自《北京航空航天大学学术论文LaTeX模板》，部分样例参考《浙江大学研究生硕士(博士)学位论文LaTeX模板》}：

\begin{tabular}{ll}
 \verb|nmu.cls |             & $\triangleright$ LaTeX宏模板文件 \\
 \verb|GBT7714-2005.bst|      & $\triangleright$ 国标参考文献BibTeX样式文件2005 \\
 \verb|GBT7714-2015.bst|      & $\triangleright$ 国标参考文献BibTeX样式文件2015 \\
 \verb|nmu_logo.png|         & $\triangleright$论文封皮北方民族大学校徽 \\
 \verb|tex/*.tex|             & $\triangleright$本模板样例中的独立章节\\
 \verb|figures/*|             & $\triangleright$本模板样例中的插图路径\\
 \verb|NMUThesis.tex|             & $\triangleright$LaTeX模板\\
 \verb|NMUThesis.pdf|             & $\triangleright$PDF模板样例\\
 \verb|make.bat|             & $\triangleright$生成NMUThesis.pdf\\
 \verb|clean.bat|             & $\triangleright$清理冗余文件\\
\end{tabular}\\

通过 \verb|\documentclass[<thesis>,<printtype>,<version>]{nmu}|载入宏包：
\begin{itemize}[leftmargin=3cm]
  \item[{\tt thesis} $\triangleright$]  论文类型(thesis)，可选值：\\
    a) 学术硕士论文（\verb|master|）[缺省值]\\
    b) 专业硕士论文（\verb|professional|）\\
    c) 博士论文（\verb|doctor|）
  \item[{\tt printtype} $\triangleright$] 打印属性(printtype)，可选值： \\
    a) 单面打印（\verb|onside|）[缺省值]\\
    b) 双面打印（\verb|twoside|）
  \item[{\tt version} $\triangleright$] 论文版本(version)，可选值： \\
    a) 盲审版（\verb|blind|）[缺省值]\\
    b) 最终版（\verb|ultimate|）
\end{itemize}

模板已内嵌LaTeX工具包：
 {\tt ifthen}，{\tt etoolbox}，{\tt titletoc}，{\tt remreset}，{\tt remreset},
 {\tt geometry}，{\tt fancyhdr}，{\tt setspace}，{\tt caption}，{\tt float}，
 {\tt graphicx}，{\tt subfigure}，{\tt epstopdf},
 {\tt book\-tabs}，{\tt longtable}，{\tt multirow}，{\tt array}, {\tt enumitem}，
 {\tt algorithm2e}，{\tt amsmath}，{\tt amsthm}, {\tt listings},
 {\tt pifont}，{\tt color}，{\tt soul}。

模板已内嵌宏：\verb|\highlight{text}|（黄色高亮）。
 
请统一使用UTF-8编码。



%-----------------------------
\section{选项设置}

\begin{itemize}[leftmargin=3cm]
  \item[{\tt  $\backslash$refcolor} $\triangleright$]  开启/关闭引用编号颜色，包括参考文献，公式，图，表，算法等\\
  \texttt{on}：开启 [默认]\\
  \texttt{off}：关闭
  \item[{\tt $\backslash$beginright} $\triangleright$]  摘要和正文从右侧开始\\
  \texttt{on}：开启 [默认]\\
  \texttt{off}：关闭
  \item[{\tt $\backslash$emptypageword} $\triangleright$]  空白页留字
  \item[{\tt $\backslash$Listfigtab} $\triangleright$]  是否使用图标清单目录\\
  \texttt{on}：开启 [默认]\\
  \texttt{off}：关闭
\end{itemize}


%-----------------------------
\section{章节撰写}
本模板支持一下标题级别标题级别

\begin{tabular}{ll}
  \verb|\chapter{章}|              & $\triangleright$ 第一章 \\
  \verb|\chapter*{无章号章}|       & $\triangleright$ 无章号章 \\
  \verb|\chaptera{无章有目录章}|   & $\triangleright$ 无章有目录章 \\
  \verb|\summary|                  & $\triangleright$ 总结\\
  \verb|\appendix|                 & $\triangleright$ 附录\\
  \verb|\acknowledgments|          & $\triangleright$ 致谢\\
  \verb|\biography|                & $\triangleright$ 个人简介\\
  \verb|\section{节}|              & $\triangleright$ 1.1 节\\
  \verb|\subsection{条}|           & $\triangleright$ 1.1.1 条\\
  \verb|\subsubsection{A}|         & $\triangleright$ 1.1.1.1 A\\
  \verb|\paragraph{a}|             & $\triangleright$ 1.1.1.1.1 a\\
  \verb|\subparagraph{a)}|         & $\triangleright$ 1.1.1.1.1.1 a)\\
\end{tabular}

%-----------------------------
\section{注意事项}
\begin{itemize}
  \item[$\triangleright$] 暂无中文斜体；
  \item[$\triangleright$] 中文粗体将转换为楷体；
  \item[$\triangleright$] 中文黑体无法加粗(盲审版封皮需单独打印)；
  \item[$\triangleright$] 行末针对标点的断行不好，例如\ref{sec:error1}处的有可能会被“、”断在了句首；
  \item[$\triangleright$] \verb|\label{<text>}|中不能使用中文；
  \item[$\triangleright$] 浮动体与正文之间的距离是弹性的；
  \item[$\triangleright$] 命令符与汉字之间请注意加空格以避免undefined错误（pdfLaTeX下好像一般不存在这个问题，主要在XeLaTeX编译环境下发生）；
\end{itemize}

%-----------------------------
\section{ToDo}
\begin{itemize}
  \item[$\triangleright$] 数学环境的行间隔；
  \item[$\triangleright$] 参考文献的行间隔；
\end{itemize}

%-----------------------------
\section{意见及问题反馈}

\indent E-mail：wizen\_zhang@163.com\\
\indent GitHub：\href{https://github.com/WizenZhang/NMUThesis/issues}{https://github.com/WizenZhang/NMUThesis/issues}


